\section*{Capítulo \ref{ch:g6:measure} --- Perímetro, área y volumen}

\begin{solution}{exo:g6:measure:1}
$3.5$ m $= 350$ cm; \quad $420$ mm $= 42$ cm; \quad
$0.8$ km $= 800$ m; \quad $25\,000$ m $= 25$ km.
\end{solution}

\begin{solution}{exo:g6:measure:2}
Rectángulo: $2 \times (8 + 3.5) = 2 \times 11.5 = 23$ cm.
Cuadrado: $4 \times 6.2 = 24.8$ cm.
Triángulo: $5 + 7 + 9 = 21$ cm.
\end{solution}

\begin{solution}{exo:g6:measure:3}
Una vuelta: $2\pi \times 50 = 100\pi \approx 314$ m. Para $2$ km
$= 2000$ m: $2000 \div 314 \approx 6.4$, así que hacen falta $7$ vueltas completas
($6$ vueltas solo hacen unos $1\,885$ m).
\end{solution}

\begin{solution}{exo:g6:measure:4}
Rectángulo: $9 \times 4 = 36$ cm$^2$. Cuadrado: $7^2 = 49$ m$^2$.
\omterm{def:g6:shapes:triangles}{Triángulo rectángulo}: $\frac{6 \times 10}{2} = 30$ cm$^2$.
Disco: $\pi \times 3^2 = 9\pi \approx 28$ cm$^2$.
\end{solution}

\begin{solution}{exo:g6:measure:5}
$3$ m$^2 = 30\,000$ cm$^2$; \quad
$45\,000$ cm$^2 = 4.5$ m$^2$; \quad
$2.5$ L $= 2\,500$ cm$^3$; \quad
$4\,500$ L $= 4.5$ m$^3$.
\end{solution}

\begin{solution}{exo:g6:measure:6}
Valla (\omterm{def:g6:measure:perimeter}{perímetro}): $2 \times (120 + 85) = 2 \times 205 = 410$ m.
\omterm{def:g6:measure:area}{Área}: $120 \times 85 = 10\,200$ m$^2$.
\end{solution}

\begin{solution}{exo:g6:measure:7}
\omterm{def:g5:solids:def}{Ortoedro}: $5 \times 4 \times 10 = 200$ cm$^3$. \omterm{def:g5:solids:def}{Cubo}:
$3 \times 3 \times 3 = 27$ cm$^3$.
\end{solution}

\begin{solution}{exo:g6:measure:8}
\omterm{def:g6:measure:perimeter}{Perímetro} $16$ cm significa largo $+$ ancho $= 8$ cm. Por ejemplo,
$6 \times 2$ (\omterm{def:g6:measure:area}{área} $12$ cm$^2$) y $5 \times 3$ (\omterm{def:g6:measure:area}{área} $15$ cm$^2$)
--- o el cuadrado $4 \times 4$ (\omterm{def:g6:measure:area}{área} $16$ cm$^2$). Cuanto más se parece a un
cuadrado, mayor es el \omterm{def:g6:measure:area}{área}.
\end{solution}

\begin{solution}{exo:g6:measure:9}
\omterm{def:g6:measure:area}{Área}: $10 \times 2 + 2 \times 6 = 20 + 12 = 32$ cm$^2$.

\omterm{def:g6:measure:perimeter}{Perímetro} (con el pie centrado bajo la barra): recorriendo el borde,
\[
10 + 2 + 4 + 6 + 2 + 6 + 4 + 2 = 36 \text{ cm}
\]
(arriba; extremo derecho de la barra; parte inferior derecha; lado derecho del pie; base del
pie; lado izquierdo del pie; parte inferior izquierda; extremo izquierdo de la barra).
\end{solution}

\begin{solution}{exo:g6:measure:10}
\emph{1.} $V = 25 \times 10 \times 2 = 500$ m$^3$.

\emph{2.} $500$ m$^3 = 500 \times 1000 = 500\,000$ L.

\emph{3.} $500\,000 \div 50\,000 = 10$ horas.
\end{solution}

\begin{solution}{exo:g6:measure:11}
Jardín: $20^2 = 400$ m$^2$. Estanque: $\pi \times 3^2 = 9\pi$ m$^2$.
Césped: $400 - 9\pi \approx 400 - 28.3 \approx 372$ m$^2$.
\end{solution}

\begin{solution}{exo:g6:measure:12}
\emph{1.} Tableta pequeña: $15 \times 6 \times 1 = 90$ cm$^3$. Tableta grande:
$30 \times 12 \times 2 = 720$ cm$^3$. El \omterm{def:g6:measure:volume}{volumen} queda multiplicado por
$720 \div 90 = 8$ --- duplicar cada una de las tres dimensiones multiplica
el \omterm{def:g6:measure:volume}{volumen} por $2 \times 2 \times 2 = 8$.

\emph{2.} Ocho veces el chocolate por cuatro veces el precio: sí, la
tableta grande sale el doble de bien por gramo.
\end{solution}

\begin{solution}{pb:g6:measure:1}
\textbf{1.} $1 \times 9$: \omterm{def:g6:measure:perimeter}{perímetro} $2 \times (1 + 9) = 20$ m,
\omterm{def:g6:measure:area}{área} $9$ m$^2$. \ $2 \times 8$: \omterm{def:g6:measure:perimeter}{perímetro} $2 \times 10 = 20$ m,
\omterm{def:g6:measure:area}{área} $16$ m$^2$.

\textbf{2.} El \omterm{def:g6:measure:perimeter}{perímetro} es el doble de (largo $+$ ancho)
(\cref{ex:g6:measure:perimeter}), así que largo $+$ ancho
$= 20 \div 2 = 10$ m. La tabla:
\begin{center}
\renewcommand{\arraystretch}{1.2}
\begin{tabular}{c|ccccc}
cercado & $1 \times 9$ & $2 \times 8$ & $3 \times 7$ & $4 \times 6$
& $5 \times 5$ \\
\hline
\omterm{def:g6:measure:area}{área} (m$^2$) & $9$ & $16$ & $21$ & $24$ & $25$ \\
\end{tabular}
\end{center}
El cuadrado $5 \times 5$ es el mejor.

\textbf{3.} $4.5 \times 5.5 = 24.75$ y $4.9 \times 5.1 =
24.99$: cada vez más cerca de $25$, pero siempre por debajo. Conjetura:
el cuadrado gana a \emph{todos} los rectángulos de \omterm{def:g6:measure:perimeter}{perímetro} $20$,
lados decimales incluidos.

\textbf{4.} De la tabla: el \omterm{def:g2:mult:def}{producto} de dos números que suman
$10$ es máximo cuando los dos números son \emph{iguales} ($5$ y
$5$). Regla general: para una \omterm{def:g1:addition:def}{suma} fija, el \omterm{def:g2:mult:def}{producto} es máximo con
partes iguales --- cuanto más desigual es la pareja, menor es el
\omterm{def:g2:mult:def}{producto}.

\textbf{5.} Del cuadrado $5 \times 5$, quita la fila de arriba
--- una tira de $5$ cuadrados unidad --- y queda un cercado de $5 \times 4$;
después pega una columna de $4$ cuadrados unidad en un extremo, y el
cercado pasa a ser $6 \times 4$. Se quitaron cinco cuadrados y solo cuatro
volvieron: la tira añadida va a lo largo del lado de $4$, que es
más corto que el lado de $5$ que cubría la tira quitada. Pérdida
neta: un metro cuadrado ($25 \to 24$). El paso siguiente, de
$6 \times 4$ a $7 \times 3$, cambia una fila de $6$ por una columna
de $3$ y pierde tres más ($24 \to 21$): cuanto más lejos está el cercado
del cuadrado, más larga es la tira que entrega y más
corta la que recibe.

\textbf{6.} Con $a + L + a = 20$, es decir, $L = 20 - 2a$:
\begin{center}
\renewcommand{\arraystretch}{1.2}
\begin{tabular}{c|ccccccc}
$w$ & $1$ & $2$ & $3$ & $4$ & $5$ & $6$ & $7$ \\
\hline
$L$ & $18$ & $16$ & $14$ & $12$ & $10$ & $8$ & $6$ \\
\omterm{def:g6:measure:area}{área} & $18$ & $32$ & $42$ & $48$ & $50$ & $48$ & $42$ \\
\end{tabular}
\end{center}
El ganador es $5 \times 10$, de \omterm{def:g6:measure:area}{área} $50$ m$^2$ --- el doble de
largo (a lo largo de la pared) que de ancho: \emph{no} es un cuadrado.

\textbf{7.} Refleja el cercado respecto de la pared: el cercado más su imagen
especular forman un cercado doble que usa la valla dos veces, $40$ m de
valla alrededor (el lado de la pared queda ahora en el interior). Por la Parte I, el
mejor rectángulo de \omterm{def:g6:measure:perimeter}{perímetro} $40$ es el cuadrado $10 \times
10$, de \omterm{def:g6:measure:area}{área} $100$ m$^2$. El mejor cercado real es la \omterm{def:g3:division:half}{mitad} de ese
mejor cercado doble: $10$ m a lo largo de la pared y $5$ m de fondo, \omterm{def:g6:measure:area}{área}
$50$ m$^2$ --- exactamente el ganador de la tabla, ahora explicado.

\textbf{8.} Longitud $20 = 2 \times \pi \times r$, así que
$r = 20 \div 6.28 \approx 3.18$ m. \omterm{def:g6:measure:area}{Área}:
$\pi r^2 \approx 3.14 \times 3.18 \times 3.18 \approx
31.8$ m$^2$ --- holgadamente más que los $25$ m$^2$ del cuadrado.
Dido sabía lo que hacía: para una longitud de borde dada, el
círculo encierra lo máximo.

\textbf{9.} \omterm{def:g6:measure:perimeter}{Perímetros}: $1 \times 36$: $74$ m; \ $2 \times 18$:
$40$ m; \ $3 \times 12$: $30$ m; \ $4 \times 9$: $26$ m; \
$6 \times 6$: $24$ m. Otra vez el cuadrado --- la menor valla para el
\omterm{def:g6:measure:area}{área} dada. Es la Parte I al revés: \omterm{def:g6:measure:perimeter}{perímetro} fijo y \omterm{def:g6:measure:area}{área} máxima,
o \omterm{def:g6:measure:area}{área} fija y \omterm{def:g6:measure:perimeter}{perímetro} mínimo: las dos coronan al cuadrado
(y, más allá de todos los rectángulos, al círculo).

\textbf{10.} Una pared redonda es el borde más corto para el espacio
que encierra (pregunta 8), así que una lata, una tubería o un depósito redondos contienen
lo máximo con el mínimo de metal --- el material es dinero.

\textbf{11.} Por ejemplo, $50$ m $\times$ $0.01$ m: \omterm{def:g6:measure:perimeter}{perímetro}
$2 \times (50 + 0.01) = 100.02$ m y \omterm{def:g6:measure:area}{área}
$50 \times 0.01 = 0.5$ m$^2$. Largo y estrecho: kilómetros de
borde y casi nada de terreno.

\textbf{12.} Falso. El cercado de la pregunta 11 tiene \omterm{def:g6:measure:perimeter}{perímetro}
$100.02$ m y \omterm{def:g6:measure:area}{área} $0.5$ m$^2$, mientras que el cercado $5 \times 5$ tiene
\omterm{def:g6:measure:perimeter}{perímetro} $20$ m y \omterm{def:g6:measure:area}{área} $25$ m$^2$: mayor \omterm{def:g6:measure:perimeter}{perímetro} y \omterm{def:g6:measure:area}{área} mucho
menor.

\textbf{13.} \omterm{def:g6:measure:area}{Área}: falta un cuadrado unidad,
$24 - 1 = 23$ cm$^2$. \omterm{def:g6:measure:perimeter}{Perímetro}: al recorrer el borde, la muesca
sustituye $1$ cm de pared recta por tres lados del cuadradito,
$1 + 1 + 1 = 3$ cm: el \omterm{def:g6:measure:perimeter}{perímetro} crece de $20$ cm a
$20 - 1 + 3 = 22$ cm. Quitar material \emph{alargó} el
borde.

\textbf{14.} Cada ampliación de una costa real revela bahías, rocas
y calas nuevas --- muescas sobre muescas ---, y la pregunta 13 muestra que
cada muesca añade longitud de borde sin apenas cambiar el \omterm{def:g6:measure:area}{área}.
Así que la longitud medida sigue creciendo con el nivel de detalle
(«la paradoja de la costa»), mientras que el \omterm{def:g6:measure:area}{área} medida se
estabiliza: el \omterm{def:g6:measure:perimeter}{perímetro} y el \omterm{def:g6:measure:area}{área} son magnitudes realmente independientes.

\textbf{15.} Argumento del espejo: el cercado doble usaría
$2 \times 36 = 72$ m de valla, y el mejor rectángulo de
\omterm{def:g6:measure:perimeter}{perímetro} $72$ es el cuadrado $18 \times 18$. Así que el mejor cercado es de
$18$ m a lo largo de la pared y $9$ m de fondo: comprobación de la valla,
$9 + 18 + 9 = 36$ m, y \omterm{def:g6:measure:area}{área} $18 \times 9 = 162$ m$^2$. En campo
abierto, el mejor es el cuadrado $9 \times 9$, de \omterm{def:g6:measure:area}{área} $81$ m$^2$.
La pared del granero \emph{duplica} exactamente el terreno del granjero.
\end{solution}
